%% kapa: exam=1 round=29 year=2018 period=2 form=A subject=appraisal_law subject_ko=감정평가관계법규 seq=21 number=21
%% source: Q-Net 감정평가사 기출문제 (cst00309) artlSeq=5202808 file=extracted/1008/5202808/2018년 제29회 감정평가사 1차시험 기출문제/2018년 제29회 감정평가사 1차시험 기출문제/2교시 A형.pdf
%% conversion: kordoc 4.24.51 -> tools/gen1.py (auto); review=proofread-sonnet; flags=none
\begin{question}{21}
국유재산법령상 일반재산에 관한 설명으로 옳지 않은 것은?
\begin{choices}
  \item 일반재산은 대부 또는 처분할 수 있다.
  \item 중앙관서의 장은 국가의 활용계획이 없는 건물이 재산가액에 비하여 유지ㆍ보수 비용이 과다한 경우 이를 철거할 수 있다.
  \item 일반재산은 매립사업을 시행하기 위하여 그 사업의 완성을 조건으로 총괄청과 협의하여 매각을 예약할 수 있다.
  \item 일반재산을 매각하는 경우에는 대통령령으로 정하는 바에 따라 매수자에게 그 재산의 용도와 그 용도에 사용하여야 할 기간을 정하여 매각할 수 있다.
  \item 총괄청은 일반재산을 보존용재산으로 전환하여 관리할 수 없다.
\end{choices}
\end{question}
\answer{5}
